\documentclass[11pt,a4paper]{article}
\usepackage[utf8]{inputenc}
\usepackage[T1]{fontenc}
\usepackage{lmodern,amsmath,amssymb,textcomp}
\usepackage[margin=24mm]{geometry}
\usepackage{longtable,booktabs,array,enumitem}
\usepackage[hidelinks]{hyperref}
\setlength{\parindent}{0pt}
\setlength{\parskip}{6pt}
\emergencystretch=3em
\hypersetup{pdfauthor={Rachel Teo, Wenyi Wang, Hamdi Abderrahmene, Alejandro Zarzuelo Urdiales}}
\begin{document}
{\LARGE\bfseries The sharp half-turn defect gap in full no-three-in-line grid configurations\par}\medskip
{\small Rachel Teo\ensuremath{^1}, Wenyi Wang\ensuremath{^1}, Hamdi Abderrahmene\ensuremath{^1}, Alejandro Zarzuelo Urdiales\ensuremath{^2}\par}\medskip
{\small\textbf{Affiliations:} \href{https://sakana.ai/}{\ensuremath{^1} Sakana AI}; \href{https://rsihouse.ai/openmath}{\ensuremath{^2} Open Math}.\par}

{\small\href{https://github.com/alejandrozu/openmath-2026-judging}{Code and formal proofs}\par}

\subsection{Abstract}
A no-three-collinear set of 2n points in an n\ensuremath{\times}n grid is either invariant under the half-turn or has at least four points whose partners are absent. An exact ten-point 5\ensuremath{\times}5 example attains four, giving the global minimum positive one-sided defect.

\subsection{Statement and counting convention}
Let A\ensuremath{\subseteq}\{0,...,n\ensuremath{-}1\}\ensuremath{^2} contain exactly 2n points, with no three distinct points collinear. Let R(x,y)=(n\ensuremath{-}1\ensuremath{-}x,n\ensuremath{-}1\ensuremath{-}y), and define the one-sided defect D=\{p\ensuremath{\in}A:R(p)\ensuremath{\notin}A\}. The theorem is |D|=0 or |D|\ensuremath{\geq}4. It is important to specify one-sided defect: |A\ensuremath{\triangle}R(A)|=2|D|, so the sharp symmetric-difference value is eight.

The theorem addresses the historically posed near-half-turn question for full grid configurations. It does not prove that a 2n-point configuration exists for every n, nor that defect four is attainable at every order.

\subsection{Geometric proof of the gap}
No row or column can contain more than two points. Since there are n rows and 2n points, each row contains exactly two, and the same holds for every column. Consequently the sum of the centred position vectors of all points of A is zero: its centroid is the grid centre.

Cancel all opposite pairs under R, together with the central fixed point if present. The remaining unmatched vectors, namely D, still sum to zero. A single unmatched vector would have to be zero and would therefore be fixed, a contradiction. Two unmatched vectors summing to zero would be opposite and therefore paired, another contradiction.

Suppose there were three unmatched points. The remaining centrally invariant part has 2n\ensuremath{-}3 points, an odd number, so it contains the grid centre. For n\ensuremath{\geq}3 this part has at least three points and therefore also contains a nonzero opposite pair. That pair and the centre are three collinear points, contradicting the hypothesis. The cases n=0,1,2 are handled separately. This centroid/parity geometry is necessary; row/column balance alone permits a two-miss alternating rectangle and does not prove the gap.

\subsection{Exact sharp witness}
The 5\ensuremath{\times}5 witness is \{(1,0),(2,0),(0,1),(4,1),(3,2),(4,2),(0,3),(1,3),(2,4),(3,4)\}. It has ten points. Exact determinants show that every triple of distinct points is noncollinear, and direct half-turn pairing gives four unmatched points. The formal proof checks those integer conditions and transports them to the grid geometry.

Independent ordinary checking confirmed all 120 distinct triples. Exhaustive enumeration through several small grids is consistent with the theorem but is supplementary, because the centroid argument proves all n. The universal theorem and witness are in the same contribution family.

\subsection{Priority and formal status}
The Bielefeld source records defect four as a minimum known example and asks about the nearest nonsymmetric configurations. The no-positive-defect-below-four theorem addresses that subsidiary question. Its historical and folklore priority remains unresolved; failure to locate an earlier proof does not establish first discovery.

The cited source site and existing defect-four configurations supply the prior setting. The theorem gives the positive-defect gap, not existence of full no-three-in-line configurations at every grid order.

\subsection{Formal verification}
The four-source development and selected sharp-gap declarations were checked with Lean using standard kernel axioms. The explicit ten-point witness and universal gap theorem are available with pinned sources and audit records in the companion repository.

\begin{samepage}
\subsection{NoThreeSymmetry.grid\_half\_turn\_gap}
\begin{quote}\small\ttfamily theorem grid\_half\_turn\_gap \{n : \ensuremath{\mathbb{N}}\} (A : Finset (GridPoint n))\newline     (h : GridNoThree A) (hcard : A.card = 2*n) :\newline     (gridDefect A).card = 0 \ensuremath{\vee} 4 \ensuremath{\leq} (gridDefect A).card\end{quote}
{\small Source: proofs/opdp43/OpHack/NoThree/Main.lean:8.\par}

{\small Source SHA-256: f8262d41b0933681b1db8dafb3d4302471dc28650b8ca950e3e1fc01f203c3d2\par}

\end{samepage}
\begin{samepage}
\subsection{NoThreeSymmetry.sharp\_half\_turn\_gap}
\begin{quote}\small\ttfamily theorem sharp\_half\_turn\_gap :\newline     (\ensuremath{\forall} (n : \ensuremath{\mathbb{N}}) (A : Finset (GridPoint n)), GridNoThree A \ensuremath{\to} A.card = 2*n \ensuremath{\to}\newline       (gridDefect A).card = 0 \ensuremath{\vee} 4 \ensuremath{\leq} (gridDefect A).card) \ensuremath{\wedge}\newline     (\ensuremath{\exists} A : Finset (GridPoint 5), A.card = 10 \ensuremath{\wedge} GridNoThree A \ensuremath{\wedge} (gridDefect A).card = 4)\end{quote}
{\small Source: proofs/opdp43/OpHack/NoThree/Main.lean:56.\par}

{\small Source SHA-256: f8262d41b0933681b1db8dafb3d4302471dc28650b8ca950e3e1fc01f203c3d2\par}

{\small\textbf{Sources and attribution:} \href{https://wwwhomes.uni-bielefeld.de/achim/no3in/symmetry_remarks.html}{1. wwwhomes.uni-bielefeld.de / symmetry\_remarks.html}; \href{https://www.unsolvedmath.com/problems/GREEN-045}{2. www.unsolvedmath.com / GREEN-045}; \href{https://github.com/alejandrozu/openmath-2026-judging/blob/d0ed487f487fa7ec814ff705ab55249983fe8707/OpenMath-Judging/review/entries/E02-opdp43-halfturn.txt}{3. alejandrozu/openmath-2026-judging / E02-opdp43-halfturn.txt}; \href{https://github.com/alejandrozu/openmath-2026-judging}{Proof source repository}.\par}

\end{samepage}
{\small \textbf{AI assistance.} Codex assisted literature checking, editorial preparation and analysis of the supplied formal artifacts. The human authors are responsible for checking the final mathematical claims, references and attribution.\par}

\end{document}
